\documentclass[12pt, a4paper]{article}

\usepackage{fontspec}
\setmainfont{TeX Gyre Pagella}
\usepackage{geometry}
\geometry{margin=1.25in}
\usepackage{setspace}
\onehalfspacing
\usepackage{amsmath}
\usepackage{amssymb}
\usepackage{amsthm}
\usepackage{titlesec}
\titleformat{\section}{\large\bfseries}{}{0em}{}
\titleformat{\subsection}{\normalsize\bfseries}{}{0em}{}
\titleformat{\subsubsection}{\normalsize\itshape}{}{0em}{}
\usepackage{parskip}
\setlength{\parindent}{1.5em}
\setlength{\parskip}{0pt}
\usepackage{booktabs}
\usepackage{array}
\usepackage{csquotes}
\usepackage{hyperref}
\hypersetup{colorlinks=false}

\theoremstyle{plain}
\newtheorem{theorem}{Theorem}[section]
\newtheorem{proposition}[theorem]{Proposition}
\newtheorem{corollary}[theorem]{Corollary}
\theoremstyle{definition}
\newtheorem{definition}[theorem]{Definition}
\theoremstyle{remark}
\newtheorem{remark}[theorem]{Remark}

\title{\textbf{Notes as Anti-Collapse}\\[0.5em]
\large\textit{A Spherepop Theory of Externalized Continuations}}
\author{Flyxion\\[0.3em]\small Independent Researcher}
\date{}

\begin{document}

\maketitle
\thispagestyle{empty}

\begin{abstract}
\noindent
The standard theory of notes treats them as passive records: artifacts that
preserve information about past events against the fallibility of biological
memory. This monograph argues that the passive theory is structurally inverted.
Notes are not primarily retrospective; they are mechanisms for preserving access
to continuations---trajectories through possibility space that would otherwise
become fragile, expensive, or unreachable. Developed within the Spherepop
framework, which treats histories as primary and states as their compressed
residues, notes emerge as the physical realization of the four fundamental
operations: Pop, Refuse, Collapse, and Bind. A taxonomy of seven note classes is
developed---Capture, Prospective, Repair, Coordination, Refusal, Generative, and
the unexpected seventh class, Collapse Notes---each characterized by its dominant
Spherepop operator. The ephemerality of a note is shown to be a mark of success
rather than failure. Connections are drawn to the MEM\textbar8 architecture, Repair
Theory, the admissibility framework, and the reachability geometry of semantic
space. The monograph concludes by arguing that civilization is best understood not
as an accumulation of knowledge but as a distributed system for preserving
continuations across generations: a gigantic, layered, partially redundant
note infrastructure whose collapse does not make a civilization ignorant but makes
it unable to reach its own knowledge.
\end{abstract}

\newpage

\tableofcontents

\newpage

%%% ============================================================
\section{The Passive Theory and Its Failure}
%%% ============================================================

There is a theory of notes so pervasive that it almost never receives explicit
statement. On this theory, a note is a device for compensating the limitations of
biological memory. Something happens---a meeting is scheduled, a formula is
derived, a perception is formed---and because the mind cannot reliably retain it,
the note is made. The note is therefore a prosthetic memory: it externalizes
information onto a more durable substrate, where it may be retrieved when
biological recall fails.

This is the passive theory of notes. It places the note in a retrospective
relation to time. The note exists because something happened in the past that is
worth preserving. The value of a note, on this account, is proportional to the
fidelity with which it reproduces a prior state: a good note accurately records
what occurred; a poor note distorts it; an ephemeral note is a low-value note
because it preserves little and does not endure.

The passive theory is not entirely wrong. Many notes are made for the reason it
describes. But as a general account of what notes are and why they matter, it
fails---and it fails in a direction that reveals something important about the
structure of cognition, language, and civilization.

The failure becomes apparent the moment one examines not celebrated notes but
common ones. A grocery list does not record anything that happened. The items on
the list do not yet exist in the form the list anticipates: they have not been
purchased, assembled, or carried home. The list records a future act. A blueprint
describes a building not yet constructed. Source code records computations not yet
run. A theorem records consequences not yet derived. A calendar records
commitments not yet honored. A playlist records a future act of listening. A
bookmark records a future act of reading. A follow list on a social network records
a future flow of information not yet received.

None of these are memories. All of them are notes.

The passive theory cannot accommodate this class of cases without becoming
incoherent. It might be extended to say that a note records an intention rather
than an event, but this extension merely postpones the difficulty. The deeper
problem is that the passive theory places every note in a backward-looking
relationship to time. Even when the recorded content is an intention, the note is
still conceived as a trace of a prior mental event. The orientation is always
retrospective.

This gets the phenomenology of note-taking almost exactly wrong. When one makes a
to-do list, one is not primarily recording a past intention; one is constructing a
future. The note is not oriented toward what was; it is oriented toward what has
not yet occurred and what might not occur without the note's assistance.

A more adequate theory begins from a different primitive.

%%% ============================================================
\section{Continuations and the Reachability Frame}
%%% ============================================================

The key distinction comes from theoretical computer science: the difference between
a \emph{state} and a \emph{continuation}.

A state is a snapshot of a system at a moment in time. It is a summary, a
compression, a projection. A continuation is what comes next---the future execution
enabled by a present configuration, the program of further development that a
state makes possible. Most theories of information focus on states. A database
stores states. A photograph is described as capturing a state. A memory is treated
as the retention of a past state.

But what makes any stored state useful is not the state itself. It is the
continuation that state makes accessible.

Consider a Git commit. The commit is not valuable because it stores bytes. It is
valuable because it preserves a path through development space---a specific
trajectory through the possibility space of a codebase that would otherwise require
expensive reconstruction. The commit does not merely record what happened to a
file; it records where the project can go from a given point, and how it can return
there from any future position. The stored state is instrumentally useful only
because it is a node in a navigable graph of continuations.

The same analysis extends across the note family. A mathematical proof is not
valuable because it stores symbols. It is valuable because it preserves a path
through inference space---a sequence of steps, each admissible given the preceding
steps, that transforms premises into conclusions, and that any competent reasoner
can traverse, extend, branch from, or identify errors within. The proof is a
compressed reconstruction path. A photograph preserves access to a perceptual
trajectory that would otherwise disappear: the image makes re-entry possible, one
can return to what the light looked like at a specific place and time, can study
details that escaped notice in the original act of perception, can share a
perceptual position across time and distance. A map does not record what happened;
it records what can be done: which routes are traversable, which locations are
accessible from which others. A song's notation or recording preserves access to a
sonic trajectory---a sequence of organized sound events that can be re-entered,
re-experienced, and extended.

This motivates the central definition of the present work.

\begin{definition}[Note]
Let $\mathcal{C}$ denote a space of continuations and let $A$ be an agent. A
\emph{note} is any artifact $N$ such that there exists a continuation
$c \in \mathcal{C}$ for which
\[
P_A(c, t \mid N) > P_A(c, t),
\]
where $P_A(c, t)$ denotes the probability that agent $A$ can successfully
reconstruct or traverse continuation $c$ at future time $t$.
\end{definition}

This definition is deliberately broad. It immediately includes bookmarks, maps,
photographs, source code, language itself, and libraries, without privileging any
particular medium. What unifies them is the functional condition: a note increases
the reachability of at least one continuation.

The medium is entirely irrelevant to the definition. The note may be textual,
visual, auditory, symbolic, computational, social, or physical. What matters is
that it maintains the traversability of a trajectory that would otherwise become
inaccessible.

\begin{definition}[Continuation Volume]
The \emph{continuation volume} of a note $N$, written $V_C(N)$, is the cardinality
of the set of continuations $c$ for which $N$ increases reachability:
\[
V_C(N) = \bigl|\{c \in \mathcal{C} : P_A(c, t \mid N) > P_A(c, t)\}\bigr|.
\]
\end{definition}

Notes differ dramatically in continuation volume. A grocery list has a continuation
volume of approximately one: it increases the reachability of a single procurement
trajectory, and almost nothing else. A mathematical proof of a fundamental theorem
may have an enormous continuation volume, enabling derivations, generalizations,
applications, and extensions across many future contexts. A programming language
has a continuation volume that is in some sense uncountable, since it creates
access to all programs expressible within it.

\begin{definition}[Reachability Leverage]
The \emph{reachability leverage} of a note $N$, written $\Lambda(N)$, is the ratio
of future reconstruction cost avoided to the cost of note creation:
\[
\Lambda(N) = \frac{\sum_{c} \bigl[C_r(c) - C_r(c \mid N)\bigr]}{C_{\text{create}}(N)},
\]
where $C_r(c)$ is the cost of reconstructing continuation $c$ from first
principles and $C_r(c \mid N)$ is the reconstruction cost given access to $N$.
\end{definition}

Reachability leverage varies over many orders of magnitude. A ten-second bookmark
may save hours of future search: high leverage. A one-line equation may preserve
centuries of accumulated derivation: extraordinarily high leverage. A commit
message written in thirty seconds may save days of future debugging: very high
leverage. Leverage analysis predicts which notes are most worth making and
preserving, quite apart from their fidelity to past states.

%%% ============================================================
\section{Notes in the Spherepop Framework}
%%% ============================================================

The Spherepop framework holds that histories are primary and states are their
compressed residues. A state $s$ is a projection
\[
\sigma : \mathcal{H} \to S
\]
that collapses many histories $h \in \mathcal{H}$ into a single observable state,
producing what the framework calls the state illusion: the appearance that the
present configuration is self-sufficient, when it is in fact the endpoint of a
history whose fibers have been discarded.

The four fundamental Spherepop operations are Pop, Refuse, Collapse, and Bind.
\emph{Pop} creates a continuation---it steps forward from a present state into a
new trajectory. \emph{Refuse} resists collapse---it preserves distinctions that
the projection $\sigma$ would otherwise destroy. \emph{Collapse} applies
$\sigma$---it compresses a history into a state, deliberately discarding detail.
\emph{Bind} couples multiple trajectory systems---it creates dependencies between
continuations that were previously independent.

Notes, viewed through this framework, emerge as the physical realization of all
four operations. A note can preserve a history (Refuse), create a new trajectory
(Pop), compress a history into a navigable summary (Collapse), or couple multiple
agents to a shared trajectory (Bind).

\begin{proposition}[Notes as Anti-Collapse Artifacts]
Let $\sigma : \mathcal{H} \to S$ be a history-to-state projection. A note $N$
is an \emph{anti-collapse artifact} if it preserves selected fibers of $\mathcal{H}$
such that a future agent can partially reconstruct elements of $\mathcal{H}$ rather
than being restricted to $S$.
\end{proposition}

\begin{proof}
By Definition 2.1, $N$ increases the reachability of at least one continuation
$c$. Any continuation $c$ that is recoverable from $N$ but not from $S$ alone
corresponds to a fiber of $\mathcal{H}$ that $\sigma$ has collapsed but that $N$
has preserved. The existence of such a continuation establishes that $N$ resists
the projection, preserving partial access to $\mathcal{H}$ against the state
illusion.
\end{proof}

This is the precise sense in which notes are anti-collapse artifacts. They do not
prevent all collapse---they cannot and should not---but they preserve selected
fibers of history that would otherwise be irretrievably lost to the projection.

The relationship between notes and Spherepop operators is not one-to-one.
Most notes mix operators. But each class of note has a dominant operation, and
this dominant operation provides the basis for a natural taxonomy. That taxonomy
is the subject of the following chapters.

%%% ============================================================
\section{Capture Notes: The Refuse Operation}
%%% ============================================================

The most familiar class of notes are those that preserve an existing trajectory
against imminent collapse. Photographs, audio recordings, journals, logs, field
notes, and measurements all belong to this class. What unites them is their
temporal orientation: they are made because a trajectory is about to disappear,
and their purpose is to preserve enough structure that a future agent may re-enter
it.

In Spherepop terms, capture notes are primarily Refuse operations. The world
naturally tends toward collapse. Distinctions disappear. Details are forgotten.
Experiences become summaries. States erase their own histories. The capture note
performs a local refusal against this tendency. It says, in effect: \emph{do not
let this distinction disappear}. It preserves enough of the fiber of history that
re-entry remains possible.

\begin{definition}[Capture Note]
A \emph{capture note} is a note whose primary function is to preserve a
continuation that exists at the time of note creation but is expected to become
inaccessible.
\end{definition}

A photograph captures a perceptual history. At the moment of exposure, a specific
configuration of light exists at a specific location. Without the photograph, this
configuration disappears: it cannot be recovered from the state of the world at
any future time, because states do not preserve their own perceptual histories.
With the photograph, the perceptual trajectory becomes re-enterable. The image is
not a replica of an experience; it is a node in a perceptual graph, a preserved
point of re-entry into a trajectory through visual space.

An audio recording captures a sonic history. The recording does not merely store
sound; it preserves access to a temporal sequence of auditory events that can be
re-entered at any point, traversed at different speeds, analyzed spectrally, or
compared against other recordings. The value of the recording is not the storage
of vibrations but the accessibility of the auditory trajectory.

A journal or laboratory notebook captures a cognitive or procedural history. The
journal entry does not merely store thoughts; it preserves the inferential
connections between them, the sequence in which conclusions were reached, the
doubts that accompanied them, and the observations that motivated them. A future
reader---including the author's future self---can re-enter the cognitive trajectory
at any point. This is why the laboratory notebook has a different epistemic status
from the published paper: the paper presents conclusions, compressing the history
of their derivation, while the notebook preserves the history itself.

\begin{proposition}[Capture Notes and Refuse]
A capture note $N$ performs a Refuse operation with respect to the projection
$\sigma : \mathcal{H} \to S$ if and only if there exists a history fiber
$h \in \sigma^{-1}(S)$ such that $P_A(c_h, t \mid N) > P_A(c_h, t)$,
where $c_h$ is the continuation associated with history $h$.
\end{proposition}

The capturing operation introduces an interesting asymmetry between the moment of
note creation and the moment of use. At creation, the note-maker has access to the
trajectory and makes a representation of it. At use, the note-reader has access
only to the representation and must reconstruct enough of the trajectory to
continue from it. The quality of a capture note is therefore not merely its
fidelity to the original but its adequacy as a reconstruction substrate.

This has consequences for how capture notes should be evaluated. A photograph that
captures the correct person but not the context in which they appear may be a poor
capture note even at high resolution, if the lost context is what matters for
future reconstruction. A field note that captures the measurements but not the
experimental conditions under which they were taken may fail to preserve the
continuation it was intended to preserve, even if every number is accurate.

The adequacy of a capture note depends on what future continuations are
anticipated. A note-maker who knows what the note will be used for can tailor it
precisely to the continuations that must remain reachable. A note-maker who does
not know what it will be used for must preserve more fiber, at greater cost, to
ensure that the relevant continuations are available to an unknown future agent.
This is the fundamental design problem of archival practice.

%%% ============================================================
\section{Prospective Notes: The Bind Operation}
%%% ============================================================

Prospective notes are the class that the passive theory cannot accommodate at all.
They include to-do lists, calendars, blueprints, roadmaps, plans, and agendas.
What distinguishes them is that they do not preserve past trajectories; they
create access to future ones.

A calendar entry is not a record of something that happened. It is a commitment
that something shall happen, and more precisely, that the trajectory from the
present to that future event shall remain reachable. Without the calendar entry,
the future event may not be attended because the continuation from present intention
to future action has been broken by the ordinary operations of memory, distraction,
and competing demands. With the entry, the continuation is stabilized: a future
self will encounter the note at the right moment and will be able to re-enter the
trajectory from intention to execution.

\begin{definition}[Prospective Note]
A \emph{prospective note} is a note whose primary function is to preserve access
to a continuation that does not yet exist at the time of note creation but is
anticipated.
\end{definition}

In Spherepop terms, prospective notes are primarily Bind operations. They bind a
future continuation to a present intention, creating a coupling across time. A
calendar binds future action to present scheduling. A blueprint binds future
construction to present design. A roadmap binds future development to a present
strategic commitment. The note creates a temporal dependency: the future
trajectory is not merely possible but linked to the present through the
note's mediation.

\begin{proposition}[Prospective Notes and Temporal Binding]
A prospective note $N$ performs a Bind operation between present agent state
$A_t$ and future continuation $c_{t'}$ if and only if
$P_A(c_{t'}, t' \mid N, A_t) > P_A(c_{t'}, t' \mid A_t)$,
where $t' > t$ and the conditioning on $A_t$ indicates that the agent's present
state is available in both cases.
\end{proposition}

The blueprint is an especially rich case. A blueprint does not merely record an
intention to construct a building; it specifies, in a form that can be handed to
agents who were not present at the original design, the precise sequence of
operations required to realize the design. The blueprint is a compressed
construction continuation: it encodes not just the desired outcome but the
admissible paths to that outcome. A contractor who has never spoken to the
architect can follow the blueprint and traverse a trajectory that the architect
defined months or years earlier.

This reveals a dimension of prospective notes that the passive theory overlooks
entirely: they can transfer continuations between agents. A to-do list written
by one person can be executed by another. A recipe written in one century can be
followed in another. A musical score written by a composer who died before it was
performed can be realized by musicians who never knew the composer. The prospective
note decouples the definition of a continuation from its execution, and can
transmit that continuation across arbitrary gaps of time, geography, and
personnel.

Source code occupies a remarkable position in this taxonomy. It is simultaneously
a prospective note and a generative artifact. Source code records computations not
yet run: it is prospective in its orientation, specifying future operations rather
than recording past ones. But it is also executable: unlike a blueprint, which
requires human agents to interpret and realize its specifications, source code can
be traversed by a machine. The program is a continuation that can execute itself.
A compiler or interpreter transforms the prospective note into an active Pop
operation.

The to-do list illustrates the simplest case of temporal binding. Its entire
function is to maintain the reachability of a set of tasks across the interval
between their identification and their completion. Without the list, the tasks
remain possible in principle but become inaccessible in practice as competing
demands displace them from working memory. With the list, the tasks are bound to
the future: a future self will encounter them, can re-enter their associated
continuations, and can execute them or reschedule them. The list does not do the
work; it preserves the reachability of the work.

%%% ============================================================
\section{Repair Notes: Refuse and Pop Combined}
%%% ============================================================

Repair notes exist because reconstruction is expensive. They include documentation,
troubleshooting guides, commit messages, lab notebooks, error logs, and post-mortem
reports. Their purpose is to preserve the path required to recover from
interruption---to make re-entry possible after a trajectory has been broken.

\begin{definition}[Repair Note]
A \emph{repair note} is a note whose primary function is to reduce the cost of
reconstructing a continuation after it has been or is anticipated to be
interrupted.
\end{definition}

Repair notes sit directly at the intersection of Refuse and Pop. The Refuse
component preserves the history of a trajectory against collapse---the documentation
refuses the loss of procedural or operational history. The Pop component enables
re-entry: when needed, the repair note provides the means to pop back into the
interrupted trajectory from an arbitrary failure state.

\begin{definition}[Reconstruction Cost]
Let $C_r(c)$ denote the cost of reconstructing continuation $c$ from first
principles, without any note. Let $C_r(c \mid N)$ denote the reconstruction cost
given access to note $N$. The \emph{repair value} of $N$ with respect to $c$ is
\[
R_v(N, c) = C_r(c) - C_r(c \mid N).
\]
\end{definition}

\begin{theorem}[Notes as Pre-emptive Repair]
Every note $N$ with positive repair value $R_v(N, c) > 0$ for some continuation
$c$ constitutes a pre-emptive repair: it reduces expected reconstruction cost
prior to any failure occurring.
\end{theorem}

\begin{proof}
By definition, $R_v(N, c) > 0$ implies $C_r(c \mid N) < C_r(c)$. Since the note
exists before any failure occurs, the cost reduction is in force from the moment
of note creation. The note therefore functions as a repair mechanism that precedes
the event it repairs for. A reactive repair responds to loss that has already
occurred; the note, by contrast, installs the means of repair in advance of any
interruption. Pre-emptive repair is structurally distinct from reactive repair:
the former reduces expected cost regardless of whether failure occurs, while the
latter responds to realized failure.
\end{proof}

The commit message is the canonical repair note in software development. A commit
message written in thirty seconds may save days of future debugging by recording
not merely what changed but why it changed---the reasoning that motivated the
modification, the alternatives that were considered and rejected, the specific
failure mode that the change was intended to address. Without this message, a
future developer who encounters the code change must reconstruct the reasoning from
the code itself, from surrounding context, from version history, or from memory---
all expensive, all lossy, all incomplete. With the message, the reasoning is
directly accessible. The repair note allows the future developer to pop back into
the inferential trajectory that produced the change.

A troubleshooting guide is a repair note for a class of failure states. It says:
if you arrive at this configuration, the continuation from here to a functional
state runs through these steps. Each entry in the guide is a stored pop
operator: given a failure state, it provides the path back to a functioning
trajectory. The guide is not valuable because it records past experiences of
failure; it is valuable because it makes future recovery cheaper.

The lab notebook occupies a complex position. As a capture note, it records
observations and measurements. As a repair note, it records the operational
procedures, instrument settings, and experimental conditions that allow the
experiment to be reproduced---not merely described, but re-entered. A lab notebook
that records results but not methods is a poor repair note, because the missing
methods are precisely what is needed to reconstruct the trajectory from the
beginning. A lab notebook that records methods in detail reduces the reconstruction
cost of the experiment to approximately zero: a future researcher can pop back into
the trajectory at any point, reproducing the conditions and continuing from where
the original investigation stopped.

%%% ============================================================
\section{Coordination Notes: Binding Multiple Agents}
%%% ============================================================

Coordination notes are distinguished not by their content but by the number of
agents they bind. Laws, standards, protocols, contracts, specifications,
interfaces, and conventions are all coordination notes. Their purpose is to
synchronize the trajectories of multiple agents who would otherwise evolve
independently.

\begin{definition}[Coordination Note]
A \emph{coordination note} is a note whose primary function is to bind the
continuations of two or more agents, creating coupled trajectories from what
would otherwise be independent evolutions.
\end{definition}

A network protocol is a coordination note. Without the protocol, two software
systems may be technically capable of exchanging information but have no means
of doing so: their continuations are independent. The protocol defines a shared
trajectory space---a set of message formats, sequencing rules, and error
conditions---within which both systems can operate simultaneously. With the
protocol, the systems are bound: each one's behavior is a continuation that
depends on the other's.

\begin{proposition}[Coordination and Collective Reachability]
Let $A_1, \ldots, A_n$ be agents with independent continuation spaces
$\mathcal{C}_1, \ldots, \mathcal{C}_n$. A coordination note $N$ increases
collective reachability if
\[
\bigl|\mathcal{C}(A_1, \ldots, A_n \mid N)\bigr|
>
\bigl|\mathcal{C}_1\bigr| + \cdots + \bigl|\mathcal{C}_n\bigr|,
\]
where $\mathcal{C}(A_1, \ldots, A_n \mid N)$ denotes the joint continuation
space available to the agents under the coordination note.
\end{proposition}

The proposition captures the supraadditive character of coordination. When agents
are bound by a coordination note, their joint trajectory space can exceed the sum
of their individual spaces because the combination of coordinated agents can
traverse paths that no individual agent could reach alone. A contract between a
manufacturer and a supplier creates joint trajectories---products that neither
party could produce independently---that are reachable only because the
coordination note has coupled their operations.

A legal system is a civilizational-scale coordination note. Laws define the
trajectory spaces within which agents may operate, specify which continuations are
admissible and which are not, and couple the behaviors of millions of agents who
have never met. The legal system is not merely a set of prohibitions; it is a
coordination infrastructure that makes large-scale cooperation possible by defining
a shared space of admissible trajectories. A society that loses its legal
infrastructure does not merely become lawless; it loses the coordination structure
that coupled its agents' trajectories, and those trajectories immediately begin
to diverge.

The bibliography at the end of a scientific paper is a coordination note about the
inferential community. It records the prior works whose continuations the present
paper extends, cites, or builds upon. The citation graph that emerges across a
scientific literature is a large-scale bind structure: each paper is coupled to the
papers it cites and to the papers that cite it, forming a network of inferential
dependencies. A paper that is uncitable---because it is unpublished, or behind an
access wall, or written in a language the community cannot read---is a paper whose
continuation has been severed from the network. Its internal inferential
continuations may be intact, but its place in the collective trajectory space has
been destroyed.

The deepest coordination notes are natural languages. A language is not merely a
communication system; it is a coordination note that defines the trajectory space
within which an entire community can interact. To share a language is to share
access to a set of semantic, syntactic, and pragmatic continuations. Words are
coordination notes that bind speakers to shared inferential trajectories: when two
speakers share the word for a concept, they share access to the same node in
semantic space and can traverse the same downstream paths from it. The death of a
language is therefore not merely the loss of a communication medium; it is the
collapse of a coordination infrastructure that coupled the semantic trajectories
of its speakers.

%%% ============================================================
\section{Refusal Notes: The Refuse Operation in Pure Form}
%%% ============================================================

Refusal notes are the closest thing the taxonomy has to a pure Refuse operation.
Mathematical proofs, cryptographic checksums, archival copies, legal records, and
formal specifications all belong to this class. Their purpose is defensive: not to
generate new trajectories or bind existing ones, but to prevent the destruction of
distinctions that already exist.

\begin{definition}[Refusal Note]
A \emph{refusal note} is a note whose primary function is to prevent the collapse
of an existing distinction, independently of whether the preserved distinction is
expected to be used.
\end{definition}

The mathematical proof is the clearest instance. A proof does not merely claim that
a theorem is true; it preserves the path through inference space that demonstrates
why it must be true, in a form that any sufficiently trained reasoner can traverse.
The proof refuses the collapse of a mathematical result into mere assertion. Without
the proof, the theorem can only be believed; with the proof, it can be known---not
because proof and belief differ in certainty, but because the proof preserves the
inferential trajectory that connects the theorem to its premises. Future
mathematicians can re-enter the derivation, verify each step, weaken assumptions,
generalize results, or identify the precise locus of errors.

\begin{proposition}[Proofs as Refusal Notes]
A mathematical proof $P$ of theorem $T$ is a refusal note that preserves the
inferential continuation from premises $\Gamma$ to conclusion $T$:
\[
P_A(c_{\Gamma \vdash T}, t \mid P) > P_A(c_{\Gamma \vdash T}, t)
\]
for any agent $A$ with sufficient background knowledge, at any future time $t$.
\end{proposition}

The proof does not merely increase reachability at the moment of creation; it
preserves the inferential trajectory indefinitely. A theorem proved in the third
century remains provable in the twenty-first by any agent who possesses the
relevant background knowledge. The proof is a refusal note of indefinite
temporal scope.

A cryptographic checksum is a refusal note in a different domain. The checksum
preserves the distinction between an intact file and a corrupted one. Without the
checksum, a corrupted file may be indistinguishable from the original; with the
checksum, the corruption is detectable. The checksum does not preserve the contents
of the file directly; it preserves the ability to distinguish the correct state of
the file from incorrect states. It is a note about the boundary of a distinction
rather than about the content within that boundary.

Archival copies preserve the reachability of documents that might otherwise be
destroyed by physical decay, technological obsolescence, or institutional
disruption. An archive does not generate new continuations; it preserves existing
ones against time. The archival impulse is therefore a Refuse impulse at
civilizational scale: the recognition that distinctions are fragile and require
active preservation if they are to remain reachable across the intervals of time
that separate their creation from their use.

The most important consequence of the refusal category is the light it sheds on
the relationship between preservation and use. Refusal notes do not require
use in order to have value. A proof that is never consulted still preserves an
inferential trajectory. An archive that is never accessed still preserves documents
against physical decay. A backup that is never restored still ensures that
restoration remains possible. The value of a refusal note is not actualized use
but preserved potential.

This has institutional implications. An archive is valuable even when no one is
consulting it, because its value lies in the maintenance of reachability, not in
current traffic. A library that is evaluated only by circulation statistics is
being evaluated by the wrong criterion: what matters is not how many books are
checked out but how many trajectories remain accessible to future agents who might
need them. The low-circulation monograph in a research library is not a failure;
it is a refusal note for an inferential trajectory that may be needed by exactly
one researcher in exactly one future context, and that would be irrecoverable if
the note were destroyed.

%%% ============================================================
\section{Generative Notes: The Pop Operation}
%%% ============================================================

Generative notes are the most powerful class in the taxonomy. They include
scientific theories, maps, programming languages, mathematical formalisms,
ontologies, and grammars. What distinguishes them from all other note classes is
that they do not merely preserve or transmit existing continuations; they create
trajectories that did not previously exist.

\begin{definition}[Generative Note]
A \emph{generative note} is a note that creates access to continuations that were
not reachable before the note's creation, independently of whether those
continuations were previously known or anticipated.
\end{definition}

In Spherepop terms, generative notes are primarily Pop operations. They step
forward from the present state of knowledge or capability into a new trajectory
space. A good theory is not merely a record of what is known; it is a continuation
factory---a compressed structure from which new knowledge can be derived
indefinitely.

\begin{proposition}[Generative Notes and Novel Reachability]
A generative note $N$ satisfies
\[
\exists\, c \in \mathcal{C} : P_A(c, t \mid N) > 0 \text{ and } P_A(c, t) = 0,
\]
for some agent $A$ and time $t$. That is, $N$ creates access to continuations
that were previously unreachable.
\end{proposition}

A map is a generative note. It does not merely record the territory that already
exists; it creates navigational possibilities that were not previously accessible.
Before a coastline is charted, the routes along it are either inaccessible or
must be discovered at great expense by direct exploration. After charting, every
agent with access to the map can traverse any route on it without independent
exploration. The map generates navigational continuations from stored spatial
structure.

A mathematical formalism is a generative note of exceptional power. When Leibniz
and Newton developed the calculus, they did not merely record a set of techniques;
they created a notation and a set of rules that made an enormous class of
previously inaccessible problems tractable. Every theorem proved in calculus
since the seventeenth century has been a continuation generated by that original
formalism. The notation itself is a generative note: the integral sign $\int$ and
the derivative $d/dx$ are compressed pop operators that make entire domains of
inference accessible to any agent who can read them.

A programming language is a generative note of a special kind: one that is
mechanically executable. The language defines not just a space of possible
programs but a machine for traversing that space. Every program expressible in
the language is a continuation that the language makes accessible. Unlike a
mathematical formalism, which requires a human reasoner to traverse its
continuations, a programming language can delegate traversal to a computer. The
language is a generative note that creates an automated pop operator.

\begin{corollary}[Generative Notes and Continuation Volume]
For a generative note $N$,
\[
V_C(N) \gg V_C(N')
\]
for any capture, prospective, repair, or coordination note $N'$ of comparable
length, because $N$ creates continuations that did not previously exist while
the other classes only preserve or bind existing ones.
\end{corollary}

Scientific theories occupy the apex of the generative hierarchy. A theory such as
Newtonian mechanics or Darwinian evolution does not merely organize existing
observations; it creates an inferential landscape within which an indefinite number
of future observations can be interpreted, new experiments can be designed, and
new predictions can be derived. The theory is a compressed reconstruction path
for the entire domain it covers: any agent with sufficient background knowledge
can enter the theory, traverse its implications, and generate new continuations
from it indefinitely.

Ontologies are generative notes that create conceptual spaces. An ontology defines
the entities, relations, and categories within a domain, and in doing so defines
what can be thought within that domain. Before an ontology exists, many
distinctions are ineffable: they can be experienced but not articulated, sensed
but not communicated. After an ontology is created, those distinctions become
reachable---nameable, combinable, transmissible. The creation of an ontology is
therefore a Pop operation at the level of thought itself: it opens trajectories
through conceptual space that were previously inaccessible.

%%% ============================================================
\section{Collapse Notes: The Hidden Class}
%%% ============================================================

The seventh class of notes is not predicted by any theory of notes as
preservation, and its existence is only visible through the Spherepop framework.
Collapse notes are artifacts whose purpose is not to preserve continuations but to
deliberately destroy detail---to compress histories into navigable states.

\begin{definition}[Collapse Note]
A \emph{collapse note} is an artifact that performs a deliberate projection
$\sigma_N : \mathcal{H} \to S_N$, compressing a history or artifact into a more
compact representation at the cost of detail, in order to make the representation
navigable.
\end{definition}

Examples are ubiquitous: file names, titles, abstracts, indexes, summaries,
taxonomies, classifications, labels, and keywords are all collapse notes. A file
name is a collapse of a document. A title is a collapse of a book. A citation is
a collapse of a paper. An abstract is a collapse of an argument. A taxonomy is a
collapse of a domain.

The paradox of collapse notes is that they appear to be the opposite of notes---
they destroy information rather than preserving it---yet they are indispensable to
the function of every other note class. Without collapse notes, the output of
generative, capture, and repair notes becomes unnavigable. An archive without a
catalogue is an archive whose contents are unreachable not because they have been
destroyed but because there is no path from need to document. An index is a
collapse note about a book: it destroys the sequential structure of the text and
replaces it with an alphabetically ordered compression, making the book's content
navigable by topic rather than by reading order.

\begin{theorem}[Necessity of Collapse Notes]
For any sufficiently large collection of notes $\mathcal{N}$, there exists a
threshold $|\mathcal{N}| > K$ above which the reachability of individual notes
in $\mathcal{N}$ decreases in the absence of collapse notes over $\mathcal{N}$.
\end{theorem}

\begin{proof}
Search cost through an unstructured collection grows at least linearly with
collection size. For $|\mathcal{N}| > K$, exhaustive search becomes too expensive
for practical use. A collapse note---index, catalogue, taxonomy---reduces search
cost by imposing structure on the collection, making individual notes reachable
at sublinear cost. Without such structure, notes become inaccessible not through
destruction but through navigational failure. Hence collapse notes are necessary
for the preservation of reachability in large collections.
\end{proof}

The theorem explains a phenomenon that the passive theory cannot: why the growth
of note collections eventually requires the development of meta-notes---notes about
notes. Libraries require catalogues. Catalogues require classification systems.
Classification systems require ontologies. The stack of collapse notes that
organizes any large archive is not overhead; it is the navigational infrastructure
that makes the archive's contents reachable at all.

The most fundamental collapse notes are linguistic. A word is a collapse note: it
compresses a complex of experience, association, inference, and usage into a
single retrievable token. The word does not preserve the history of the experiences
it compresses; it collapses them into a navigable unit that can be combined with
other units to traverse semantic space. The cost of this compression is the loss
of whatever individual experiences were not preserved in the word's semantic
content. The benefit is the vast increase in semantic navigability that shared
vocabulary affords: speakers with common vocabulary can traverse shared semantic
trajectories far more efficiently than speakers who must describe each concept
from experiential primitives.

A noun is therefore a collapse note about a process. The passive theory of notes
would not recognize this as a note at all: nouns do not preserve information about
specific past events. But the continuation theory recognizes them immediately as
the most pervasive and powerful class of collapse notes that civilization has
developed, enabling the compression of unlimited historical complexity into
navigable semantic tokens.

%%% ============================================================
\section{The Ephemerality Inversion}
%%% ============================================================

With the taxonomy in place, the ephemerality question can be addressed precisely.
The passive theory uses persistence as a criterion of note value: a note that
survives is more valuable than one that disappears, because survival indicates
successful preservation of the past state. The continuation theory predicts the
opposite.

\begin{definition}[Continuation Completion]
A continuation $c$ is \emph{complete} with respect to note $N$ at time $t$ if
the trajectory encoded by $N$ has been successfully traversed: the agent has
arrived at the terminal state of $c$ and no further traversal is required.
\end{definition}

\begin{theorem}[Ephemerality Inversion]
Let $N$ be a note encoding a single continuation $c$. If $c$ has been completed,
then the destruction of $N$ does not reduce the reachability of any remaining
continuation. Therefore, the ephemerality of $N$ after the completion of $c$
implies neither loss nor failure.
\end{theorem}

\begin{proof}
By assumption, $c$ is the sole continuation for which $N$ increases reachability.
After completion of $c$, no continuation remains for which $P_A(c', t \mid N)
> P_A(c', t)$ for any $c' \neq c$. Therefore $V_C(N) = 0$ at time $t$.
Since continuation volume is zero, the destruction of $N$ reduces reachability
by zero. The note has fulfilled its function, and its destruction represents
the completion of a successful operation, not a loss.
\end{proof}

The grocery list that is discarded at the supermarket exit has not failed. It has
succeeded completely. The continuation it encoded---the trajectory from home to
supermarket to provisioned kitchen---has been traversed. The note's continuation
volume has been reduced to zero by the execution of its purpose. Its
disappearance is the correct outcome of a successful note, not a mark of
low quality.

The theorem generates a classification of possible note fates that inverts the
passive theory's ranking.

\begin{center}
\begin{tabular}{lll}
\toprule
Fate & Passive Theory Ranking & Continuation Theory Ranking \\
\midrule
Executed and discarded & Low & Optimal \\
Preserved and re-entered & Medium & Valuable \\
Preserved, never re-entered & High & Unrealized potential \\
\bottomrule
\end{tabular}
\end{center}

The passive theory ranks preservation as inherently good, so the note that
survives longest is ranked highest. The continuation theory ranks the completion
of function as good: the note that is used and discarded has realized its value
fully, the note that is used and preserved continues to offer value, and the note
that is preserved but never used has stored value that was never realized.

This has consequences for how archivists, librarians, and information managers
should evaluate their collections. A collection of notes with high usage and high
discard rates is performing well: continuations are being executed. A collection
with high preservation and low usage may be failing: it has stored potential that
is not being actualized. The relevant question is not how many notes survive but
how many continuations remain reachable and are being traversed.

The historical record tends to preserve notes whose continuations were not
completed. The unfinished symphony survives because it was never finished. The
ancient letter survives because the correspondence it initiated was interrupted.
The architectural blueprint of a demolished building may outlast the building
itself. The most celebrated survivors of antiquity---incomplete manuscripts,
undelivered letters, abandoned plans---are celebrated in part because their
incompleteness makes their original continuations visible. The completed notes,
the sent letters, the finished and demolished buildings, have left no trace
precisely because they succeeded.

%%% ============================================================
\section{Why LaTeX Preserves More}
%%% ============================================================

The continuation volume framework resolves a question that matters for anyone who
has switched from prose writing to mathematical writing and noticed that the
experience is different in a way that is difficult to articulate. The difference
is not aesthetic, and it is not about typographical quality---though the typography
is genuinely better. The difference is that different note substrates preserve
different classes of continuation.

\begin{proposition}[Substrate-Specific Continuation Classes]
Different note substrates preserve fundamentally different classes of continuation:
\begin{itemize}
\item A plain prose document preserves linguistic continuations.
\item A musical score preserves performance continuations.
\item Source code preserves computational continuations.
\item A \LaTeX\ document preserves mathematical and inferential continuations.
\item A Git repository preserves developmental continuations.
\item A map preserves navigational continuations.
\end{itemize}
For each substrate, continuations of the preserved class are accessible to agents
with appropriate reading competencies; continuations of other classes are not
preserved in their full form.
\end{proposition}

A plain text paraphrase of a theorem preserves approximately zero inferential
continuations from the theorem. It allows a reader to know that a certain result
holds but not to verify it, extend it, apply it, or identify its assumptions. The
paraphrase has collapsed the theorem into a bare assertion, discarding the proof
and all the inferential structure that gives the theorem its technical content.

A \LaTeX\ document containing the theorem and its proof preserves the full
inferential trajectory. A reader can enter the proof at any step, verify the
inference, identify the lemmas being used, trace the assumptions backward to their
sources, and extend the argument to new cases. The proof is a stored path through
inference space: it does not merely assert the destination but records the route,
in a form that can be traversed by any agent with sufficient background knowledge.

\begin{proposition}[LaTeX and Continuation Volume]
For a theorem $T$ with proof $P$ and prose paraphrase $D$ of comparable length,
\[
V_C(P) \gg V_C(D),
\]
because $P$ preserves the inferential continuations from premises to conclusion
while $D$ preserves only the linguistic continuation of the assertion.
\end{proposition}

The claim extends beyond theorems. A \LaTeX\ monograph that contains equations,
commutative diagrams, defined environments, and cross-referenced propositions
preserves a richer continuation space than a prose document of the same length,
because the formal environments are nodes in a navigable inferential graph. Each
defined term is a node. Each theorem is a path. Each proof is a traversal of that
path, recorded in a form that can be replayed.

The reachability leverage of a mathematical monograph is correspondingly higher
than that of a prose essay. The monograph compresses an extensive inferential
landscape into a few hundred pages: thousands of person-years of mathematical
development, in a form that any competent reader can traverse in hours. This is
the deep reason why mathematical literature accumulates differently from prose
literature: each mathematical result, once proved and published, is available as a
continuation for all future mathematical work, without diminution, without requiring
re-derivation, accessible to anyone who can read the notation. The mathematical
literature is a gigantic generative and refusal note structure, and \LaTeX\ is
the substrate that has made its continuations most accessible.

%%% ============================================================
\section{Connection to MEM\textbar8}
%%% ============================================================

The MEM\textbar8 architecture treats memory not as isolated symbol storage but as
structured pattern access, where retrieval depends on preserving relationships,
locality, and accessibility. The storage layer is a cognitive substrate rather than
a passive database, with locality elevated to a principle of cognitive coherence.

Viewed through the note taxonomy, MEM\textbar8 is an architecture for managing the
full range of note classes within a single cognitive system. Each layer of
MEM\textbar8 corresponds to a different note class operating at a different
temporal and spatial scale.

\begin{proposition}[Notes and MEM\textbar8 Layers]
The note taxonomy maps onto MEM\textbar8 architecture as follows. Capture notes
correspond to the episodic layer: they preserve specific events in their temporal
and contextual specificity, enabling later reconstruction of the associated
trajectories. Repair notes correspond to the procedural layer: they record
operational sequences that reduce reconstruction cost when skills or routines are
interrupted. Coordination notes correspond to the semantic layer: shared
conceptual structures that couple individual cognitive trajectories to community
norms. Collapse notes correspond to the indexical layer: compressed access
structures that make the contents of other layers navigable. Generative notes
correspond to the theoretical layer: frameworks that create new inferential
trajectories from stored patterns.
\end{proposition}

The MEM\textbar8 principle that memory is event-structured rather than symbol-
structured aligns precisely with the continuation theory of notes. A memory is not
a stored state; it is a stored path---a trajectory through event space that can be
re-entered and traversed. The retrieval process is not lookup but reconstruction:
the agent does not retrieve a stored item from a fixed address but reconstructs a
trajectory from available cues, following the path back through the event structure
to the desired content.

This event-structured view of memory implies that forgetting is not the loss of
stored symbols but the collapse of trajectories. What is forgotten is not a datum
but a path: the sequence of connections through event space that made a particular
memory reachable. Notes, on this view, are external trajectory preservers: they
maintain paths through event space that the internal event-structure has lost or
is at risk of losing.

A notebook is a miniature MEM\textbar8. A Git repository is a miniature
MEM\textbar8. An archive is a miniature MEM\textbar8. Each stores not symbols but
structured trajectories, organized in a way that makes reconstruction possible.
The difference between a good archive and a poor one is not primarily the quantity
of material stored but the quality of the access structure: how well the collapse
notes and coordination notes that organize the collection preserve the
reachability of the trajectories within it.

%%% ============================================================
\section{Connection to Repair Theory}
%%% ============================================================

Repair Theory holds that repair is the restoration of access to a trajectory that
has been interrupted. A broken tool is one whose operational continuations have
been disrupted; to repair it is to restore the trajectory from intention to
execution. A corrupted file is one whose informational continuations have been
lost; to recover it is to restore a path through the space of derivable states.

Notes are pre-emptive repairs, as established in Theorem 6.2. But the connection
runs deeper. Repair Theory implies a classification of failure modes that maps
directly onto the note taxonomy.

Each note class is a pre-emptive response to a specific class of anticipated
failure. Capture notes are pre-emptive repairs for perceptual and experiential
collapse: the photograph is made in anticipation of the failure of biological
visual memory. Prospective notes are pre-emptive repairs for intention-execution
gaps: the calendar entry is made in anticipation of the failure to maintain
continuity between a scheduled intention and a future moment of execution. Repair
notes proper are pre-emptive repairs for procedural interruption. Coordination
notes are pre-emptive repairs for coordination failure: the protocol is established
in anticipation of the failure that would occur if agents evolved their behaviors
independently. Refusal notes are pre-emptive repairs for distinction collapse.
Generative notes are pre-emptive repairs of a different kind: they create
trajectory spaces in anticipation of future needs that may not yet be articulable.

\begin{theorem}[Note Taxonomy as Failure Taxonomy]
The seven classes of notes in the taxonomy correspond bijectively to seven classes
of anticipated trajectory failure: experiential collapse, intention-execution gap,
procedural interruption, coordination failure, distinction collapse, navigational
failure, and inferential poverty.
\end{theorem}

The theorem implies that the presence of a particular note class in a community
or institution is evidence of the particular failure modes that community
anticipates. A community that develops extensive archival practices anticipates
distinction collapse. A community that develops extensive documentation practices
anticipates procedural interruption. A community that develops extensive legal
systems anticipates coordination failure.

The civilization question becomes, in this framework, a question about which
failure modes a civilization is equipped to anticipate and address. A civilization
with a rich note infrastructure is one that has identified the failure modes most
threatening to its trajectory spaces and has developed specific note technologies
to address each of them. The gaps in a civilization's note infrastructure---the
failure modes it has not anticipated---are the loci of its most catastrophic
vulnerabilities.

%%% ============================================================
\section{Civilization as a Note System}
%%% ============================================================

The argument of the preceding chapters converges on a claim about civilization.
The standard narrative describes civilization as the accumulation of knowledge:
civilizations advance by discovering more, learning more, and recording more. The
continuation theory suggests a different description: civilization is the
construction and maintenance of a distributed note infrastructure for preserving
trajectories across time.

\begin{definition}[Civilizational Note Infrastructure]
Let civilization $Z$ possess note infrastructure $\mathcal{N} = \{N_1, \ldots, N_k\}$.
Define collective reachability as
\[
R(Z) = \bigcup_{i=1}^{k} \mathcal{C}(N_i),
\]
where $\mathcal{C}(N_i)$ is the set of continuations preserved or enabled by $N_i$.
\end{definition}

\begin{theorem}[Civilizational Collapse]
The destruction of note infrastructure $\mathcal{N}$ decreases collective
reachability even when the underlying facts, truths, or states that the notes
describe remain unchanged:
\[
R(Z \setminus \mathcal{N}) < R(Z).
\]
\end{theorem}

\begin{proof}
By Definition 13.1, $R(Z)$ includes all continuations preserved by elements of
$\mathcal{N}$. If $\mathcal{N}$ is destroyed, no element of $\mathcal{N}$
contributes to collective reachability. Since some continuations are accessible
only through $\mathcal{N}$---that is, $\mathcal{C}(N_i) \not\subseteq R(Z
\setminus \mathcal{N})$ for some $i$---it follows that $R(Z \setminus \mathcal{N})
< R(Z)$. The facts or truths themselves are not destroyed, but the paths that
made them accessible are. A civilization that cannot reach its truths is
functionally indistinguishable from one that does not possess them.
\end{proof}

The theorem formalizes an observation that is easy to state but hard to fully
absorb: a civilization that loses its note infrastructure does not become ignorant.
It becomes unable to reach its own knowledge. The tragedy of lost languages, burned
libraries, undocumented codebases, broken hyperlinks, dead file formats, and
uncatalogued archives is not that information ceased to exist. It is that
continuations ceased to be reachable.

The progression of note technologies through history is intelligible as a sequence
of expansions in civilizational reachability. Speech allowed continuations to
survive the duration of a conversation. Stories allowed them to survive a lifetime.
Writing allowed them to survive centuries. Printing allowed them to survive across
populations that share no living memory. Digital storage allowed them to survive
replication at essentially zero marginal cost. Version control allowed them to
survive modification, preserving access not just to what a document is now but to
what it was at every prior point and how it got here. The internet allows
continuations to survive distribution across arbitrary geography.

\begin{center}
\begin{tabular}{lll}
\toprule
Technology & Reachability Horizon & Dominant Note Classes \\
\midrule
Speech & Minutes to hours & Coordination, Capture \\
Story & Decades & Capture, Generative \\
Writing & Centuries & All classes \\
Printing & Populations & Coordination, Refusal \\
Digital storage & Replication & All classes, at scale \\
Version control & Modification history & Repair, Capture \\
Network & Distribution & Coordination, Generative \\
\bottomrule
\end{tabular}
\end{center}

Each note technology expands the reachability horizon of civilizational
continuations by addressing a specific class of failure that the preceding
technology could not prevent. Writing prevents the death of witnesses. Printing
prevents the destruction of individual copies. Digital storage prevents the cost
barriers to copying. Version control prevents the loss of developmental history
through modification. Each technology is a response to the failure modes that the
preceding technology left unaddressed.

Libraries are not repositories of knowledge in the passive sense. They are
repositories of continuations---systems for preserving the traversability of
trajectories across the interval between a document's creation and a reader's
need. The catalogue is a collapse note about the collection: it preserves access
to the notes themselves by encoding their location within the system. An
uncatalogued archive is an archive whose outer reachability structure has failed,
even if the inner notes remain physically intact.

Scientific journals are note repositories that have developed quality-control
mechanisms for the notes they store. Peer review is not primarily a system for
certifying truth; it is a system for verifying that the continuations a paper
encodes are genuine---that the steps are admissible, that the derivations hold,
that the experiments can be reproduced. A paper that fails peer review has not
recorded the wrong past; it has published broken continuations, paths that claim
to reach destinations they do not in fact reach.

%%% ============================================================
\section{The Admissibility Connection}
%%% ============================================================

The admissibility framework holds that a trajectory is admissible relative to a
set of constraints if it satisfies those constraints at every point of its
execution. The intersection of notes with admissibility theory is deep: notes do
not merely preserve continuations, they preserve \emph{admissible} continuations,
and the failure of a note may consist not in the destruction of a continuation but
in the collapse of its admissibility structure.

A proof is a refusal note for an inferential trajectory. But it is more
specifically a refusal note for the admissibility of that trajectory: it preserves
not just the conclusion but the evidence that each step in the derivation is
admissible---that the inference rules used are valid, that the assumptions invoked
are available, that the quantifiers are properly scoped. A proof that omits a step
is a proof with a gap in its admissibility record: the trajectory may still be
traversable by an expert who can fill the gap, but the note has failed to preserve
the admissibility verification for future agents who cannot.

\begin{proposition}[Notes and Admissibility Preservation]
A note $N$ preserves an admissible continuation $c$ relative to constraint set
$\Omega$ if and only if $N$ encodes sufficient information to verify the
admissibility of $c$ with respect to $\Omega$ at each step, independently of
any information not contained in $N$.
\end{proposition}

This proposition reveals a dimension of note quality that goes beyond fidelity or
continuation volume. A note may increase the reachability of a continuation while
failing to preserve its admissibility structure. A traveller's account of a route
may enable future travellers to attempt the route while failing to record which
segments are only traversable in dry weather: the continuation is preserved but
its admissibility conditions are not. The note has preserved the trajectory while
losing the admissibility geometry.

The xylomorphic criterion from the admissibility framework---the condition that a
trajectory must satisfy $\lambda < 1$ to avoid divergence---implies a class of
notes that might be called admissibility certificates. An admissibility certificate
is a note that preserves not the trajectory itself but the evidence of its
admissibility: a test result, a formal verification, a peer review report. The
certificate decouples the question of whether a trajectory can be traversed from
the question of whether it is safe or valid to traverse it.

Version control systems maintain admissibility certificates in a distributed form.
The commit history is not merely a record of developmental states; it is a record
of the decisions made at each point, the tests that passed, the reviews that were
conducted. This admissibility record is what makes the history navigable not just
for reconstruction but for evaluation: a future developer can not only return to
any prior state but assess the conditions under which that state was considered
admissible.

%%% ============================================================
\section{Conclusion: The Algebra of Continuations}
%%% ============================================================

The argument of this monograph can be compressed into a small number of
propositions.

Notes are not passive records of the past. They are anti-collapse artifacts:
mechanisms that preserve selected fibers of history against the projection that
reduces histories to states. A note is any artifact that increases the reachability
of at least one continuation that would otherwise become inaccessible.

The standard evaluation criterion for notes---persistence---is the wrong criterion.
The correct criterion is continuation realization. A note that is used and discarded
has realized its value fully. A note that is preserved indefinitely but never
consulted has stored value that was never actualized. The ephemerality of a note
is a mark of success when the continuation it encoded has been executed.

The Spherepop framework reveals that notes are the physical realization of the four
fundamental operations. Refuse gives capture notes and refusal notes. Pop gives
generative notes. Bind gives prospective notes and coordination notes. And Collapse
gives the seventh and most counterintuitive class: collapse notes, the artifacts
whose purpose is controlled forgetting---the file name, the title, the index, the
taxonomy---without which the output of all other note classes becomes unnavigable.

The civilization that loses its note infrastructure does not become ignorant. It
becomes unable to reach its own knowledge. The historical record of civilizational
collapse, cultural loss, and technological forgetting is, at the level of mechanism,
a record of note infrastructure failure: the collapse of the trajectories that made
knowledge accessible, rather than the destruction of knowledge itself.

And the history of civilizational progress is, at the same level of mechanism,
a history of expanding note technologies. Each major development---writing,
printing, the scientific journal, the archive, the digital repository, version
control---is the construction of a new class of note infrastructure that addresses
a failure mode the preceding infrastructure could not prevent. The work of
civilization is largely the work of extending the reachability horizon of a
species across time.

\medskip

Spherepop is the algebra of continuations. Notes are the physical realization of
that algebra. And civilization is the attempt, always incomplete, always under
threat, to build a note infrastructure adequate to the continuation needs of the
species that requires it.

\newpage

\begin{thebibliography}{99}

\bibitem{spherepop}
Flyxion. \emph{Spherepop: A Calculus of Collapse and Continuation}. Monograph.
Independent Researcher.

\bibitem{programs_as_histories}
Flyxion. \emph{Programs as Histories: The State Illusion and Its Consequences}.
Monograph. Independent Researcher.

\bibitem{mem8}
Flyxion. \emph{MEM\textbar8: Memory as Event Structure}. Monograph. Independent
Researcher.

\bibitem{repair_theory}
Flyxion. \emph{Repair as a Fundamental Category}. Monograph. Independent
Researcher.

\bibitem{reachability}
Flyxion. \emph{From Preservation to Reachability: Semantic Continuity in an Age
of Drift}. Monograph. Independent Researcher.

\bibitem{admissibility}
Flyxion. \emph{The Admissibility Field}. Monograph. Independent Researcher.

\bibitem{strachey}
Strachey, C. Fundamental concepts in programming languages. \emph{Higher-Order
and Symbolic Computation} 13 (2000): 11--49. Originally delivered 1967.

\bibitem{felleisen}
Felleisen, M. and Friedman, D.~P. A calculus for assignments in higher-order
languages. \emph{Proceedings of the 14th ACM SIGACT-SIGPLAN Symposium on
Principles of Programming Languages}, 1987.

\bibitem{hutchins}
Hutchins, E. \emph{Cognition in the Wild}. MIT Press, 1995.

\bibitem{clark}
Clark, A. and Chalmers, D. The extended mind. \emph{Analysis} 58, no.~1 (1998):
7--19.

\bibitem{goody}
Goody, J. \emph{The Domestication of the Savage Mind}. Cambridge University
Press, 1977.

\bibitem{ong}
Ong, W.~J. \emph{Orality and Literacy: The Technologizing of the Word}. Methuen,
1982.

\bibitem{eisenstein}
Eisenstein, E.~L. \emph{The Printing Press as an Agent of Change}. Cambridge
University Press, 1980.

\bibitem{borgman}
Borgman, C.~L. \emph{Scholarship in the Digital Age: Information, Infrastructure,
and the Internet}. MIT Press, 2007.

\bibitem{atkinson}
Atkinson, R.~C. and Shiffrin, R.~M. Human memory: A proposed system and its
control processes. \emph{The Psychology of Learning and Motivation} 2 (1968):
89--195.

\bibitem{tulving}
Tulving, E. Episodic and semantic memory. In E. Tulving and W. Donaldson (eds.),
\emph{Organization of Memory}. Academic Press, 1972.

\end{thebibliography}

\end{document}
